% Chapter introduction --- plan \S4.1 (Phase 5).
% Absorbs the killing/catastrophe distinction from the former
% 02_specification_of_rupture_state.tex (typos fixed).

\section{Introduction}

The standard models of within-host infection treat the infected cell as a
black box. In the Basic Model of Viral Replication, a cell that has been
infected releases new particles at a constant rate $p$ until it is removed
at a constant rate $d_{\Icell}$; nothing about the cell's interior enters
the dynamics, and the two constants are phenomenological --- parameters to
be fitted, not quantities to be derived. For a budding virus this is at
least a defensible caricature. For a lytic pathogen it is plainly wrong in
mechanism, whatever it may achieve by fitting: such a pathogen grows inside
the cell, and releases nothing until the cell ruptures, at which point it
releases everything at once. \textit{Yersinia pestis} in the macrophage,
\textit{Francisella tularensis} in the macrophage, the lytic phage in its
bacterium --- all follow the intracellular replication-and-lysis picture,
and all are poorly served by a constant drip.

This chapter asks what the constant-rate assumption should be replaced
with, and what the replacement implies. The intracellular model was set up
in Chapter~3: the birth--death--catastrophe process, in which the
pathogen count within one cell is born at rate $\beta$ per particle, dies
at rate $\mu$ per particle, and triggers the rupture of the cell at rate
$\delta$ per particle. Chapter~3 derived the principal single-cell
quantities --- the no-rupture probability $I(t)$, the mean and second
moment $J(t)$, $K(t)$, and the cumulative mean release $V(t)$. What it did
not derive is the full distribution of the process, the statistics of the
rupture event itself, or the passage from one cell to a population of them.
These are the subjects of the present chapter.

Four questions organise it.
\begin{enumerate}
\item What is the complete single-cell description? The state
probabilities at every time, the quasi-stationary distribution of the
load, and the laws of the burst time and the burst size
(\S\S\ref{sec:state-probs}--\ref{sec:cond-means}).
\item How do several founders in one cell, and chains of cells, behave?
Multiplicity of infection (\S\ref{sec:moi}) and chained immediate transfer
(\S\ref{sec:chained}).
\item How should the single-cell theory be embedded in a population model?
Why the constant-rate substitution fails (\S\ref{sec:why-constant-fails}),
and the age-structured renewal model that replaces it
(\S\ref{sec:renewal}).
\item What does the replacement imply? Effective parameters and the
invariance of $R_0$, and what population data can and cannot identify
(\S\ref{sec:peff-section}); the flooding advantage of bursting and its
trade-off against invasion speed (\S\ref{sec:flooding}); the spectrum of
release models between bursting and budding (\S\ref{sec:spectrum}); and the
boundary of the lytic claims, drawn against HIV (\S\ref{sec:hiv}).
\end{enumerate}

The answers, in advance and in brief, are these. The single-cell process
is solvable far beyond its moments: the state probabilities are geometric
at every time, and their conditional limit --- the quasi-stationary
distribution --- is geometric with ratio $1/a$. The burst-size distribution,
conditioned on bursting, is \emph{exactly the same law}: a cell's rupture
size is distributed as the load it would have had, had it kept growing
forever. A population of such cells is not a constant-rate system but an
age-structured one, governed by two closed-form kernels --- the survival
kernel $\Ihat(a)$ and the release kernel $g(a)=\delta K(a)$ --- and the
renewal model built on them contains classical BMVR as an exact
exponential-phase projection. Under that projection, the effective release
rate $p_{\rm eff}$ depends on the epidemic growth rate, falling by orders
of magnitude from the stationary phase to the acute phase; $R_0$ itself is
invariant to bursting at fixed total output; and the three intracellular
rates cannot be recovered from the two BMVR parameters alone. At fixed
output, bursting lowers the probability of early extinction --- relative to
matched budding --- if and only if $1/\delta<1/\beta+1/\mu$, at the cost of
a slower deterministic invasion. And the honest boundary of all this is
that HIV, which buds, is not a BDC pathogen; the chapter ends by saying
precisely what the lytic theory does and does not offer it.

\subsection{Two varieties of rupture: killing and catastrophe}
\label{sec:rupture-conventions}

Before we begin, it is important to state a distinction between two
varieties of the model. When a rupture event occurs, the process goes to a
state. In the literature, this state is sometimes chosen to be $0$, and in
other cases there is a separate rupture state defined, sometimes labelled
$R$. The analysis does differ between the two options. And so for clarity,
in this chapter we will refer to the $0$-state case as ``killing''; the
distinct rupture state as ``catastrophe''. The distinction matters most in
\S\S\ref{sec:pgf-pdes}--\ref{sec:state-probs}, where both conventions are
carried in parallel; from \S\ref{sec:qsd} onward the catastrophe
convention is primary.
